\subsection*{4.2 Divergence of a Vector Field}

\begin{conceptbox}
    The divergence measures how much a vector field spreads out (source) or converges (sink):
    \begin{equation*}
        \nabla \cdot \mathbf{F} = \frac{\partial P}{\partial x} + \frac{\partial Q}{\partial y} + \frac{\partial R}{\partial z}
    \end{equation*}
\end{conceptbox}

\begin{examplebox}
    Find the divergence of $\mathbf{F}=\langle x,y,z\rangle$.

    \textbf{Step 1: Identify components}
    \begin{align*}
        P=x,\quad Q=y,\quad R=z
    \end{align*}

    \textbf{Step 2: Compute partial derivatives}
    \begin{align*}
        \frac{\partial P}{\partial x} = 1,\quad
        \frac{\partial Q}{\partial y} = 1,\quad
        \frac{\partial R}{\partial z} = 1
    \end{align*}

    \textbf{Step 3: Sum them}
    \begin{align*}
        \nabla \cdot \mathbf{F} = 1+1+1 = 3
    \end{align*}

    \textbf{Interpretation:} The field has a constant positive divergence, meaning it behaves like a uniform source everywhere.
\end{examplebox}

\begin{examplebox}
    Find divergence of $\mathbf{F}=\langle x^2, y^2, z^2 \rangle$.

    \begin{align*}
        \nabla \cdot \mathbf{F}
         & = \frac{\partial}{\partial x}(x^2) + \frac{\partial}{\partial y}(y^2) + \frac{\partial}{\partial z}(z^2) \\
         & = 2x + 2y + 2z
    \end{align*}

    \textbf{Interpretation:} Divergence varies with position; stronger away from origin.
\end{examplebox}